% !TeX program = pdflatex
% IEEE Access adaptation of the current main-SMPT manuscript.
% The class reuses TeX's \year name; PGF needs the original numeric primitive.
\let\TeXyear\year
\documentclass{ieeeaccess}
\let\AccessYear\year
\let\year\TeXyear
\usepackage[T1]{fontenc}
\usepackage[utf8]{inputenc}
\usepackage[english]{babel}
\usepackage{cite}
\usepackage{amsmath,amssymb,amsfonts}
\usepackage{graphicx}
% Permit wide figures at the bottom as well as the top of a text page.
\usepackage{stfloats}
\usepackage{tikz}
\usepackage{pgfplots}
\usepgfplotslibrary{groupplots,fillbetween}
\usetikzlibrary{patterns,positioning,fit,arrows.meta}
\usepackage{booktabs,array,float,xcolor}
% Review replacements: new text is royal blue; original text is retained.
\definecolor{RoyalBlue}{RGB}{65,105,225}
\newcommand{\michael}[2]{{\color{RoyalBlue}#1}}
% Yulian uses the same new/previous-text convention, with changes in red.
\newcommand{\yulian}[2]{{\color{red}#1}}
% Use the template blue as process CMYK; PGF also manages PDF colour resources.
\definecolor{accessblue}{cmyk}{1,0.3,0,0.2}
\definecolor{greycolor}{cmyk}{0,0,0,0.8}
\definecolor{grey}{cmyk}{0,0,0,0.1}
\usepackage{microtype}
\microtypesetup{expansion=false}
\IfFileExists{xurl.sty}{\usepackage{xurl}}{\usepackage{url}}
\urlstyle{same}
\DeclareUrlCommand\path{\urlstyle{same}}
\setlength{\emergencystretch}{2em}
\hyphenation{inter-community intra-community Gillespie}
\floatstyle{ruled}
\newfloat{algorithm}{tbp}{loa}
\floatname{algorithm}{Algorithm}
\newtheorem{proposition}{Proposition}
\newtheorem{corollary}{Corollary}
\newtheorem{remark}{Remark}
\newenvironment{proof}{\begin{IEEEproof}}{\end{IEEEproof}}
\newcommand{\ibar}{\bar{y}}
\newcommand{\sbar}{\bar{s}}
\newcommand{\Exp}{\mathrm{Exp}}
\newcommand{\Full}{\textup{Full}}
\newcommand{\TS}{\textup{TwoScale}}
\newcommand{\ind}{\mathbf{1}}
\newcommand{\eW}{e_{W}}
% Keep the supplied footer's appearance without its overfull zero-width boxes.
\makeatletter
\newcommand{\AccessOddFooter}{\hbox to\textwidth{{\footervolfont VOLUME\ \thevol, \theyear}\hfil{\footerpagefont\thepage}}}
\newcommand{\AccessEvenFooter}{\hbox to\textwidth{{\footerpagefont\thepage}\hfil{\footervolfont VOLUME\ \thevol, \theyear}}}
\let\AccessTitleStyle\ps@titlepage
\let\AccessHeadingStyle\ps@headings
\def\ps@titlepage{\AccessTitleStyle
  \def\@oddfoot{\raisebox{3pt}{\AccessOddFooter}}
  \def\@evenfoot{\raisebox{3pt}{\AccessEvenFooter}}}
\def\ps@headings{\AccessHeadingStyle
  \def\@oddfoot{\raisebox{8pt}{\AccessOddFooter}}
  \def\@evenfoot{\raisebox{8pt}{\AccessEvenFooter}}}
\makeatother
\pagestyle{headings}
% Keep float spacing compact while retaining the IEEE page dimensions.
\setlength{\dbltextfloatsep}{12pt plus 2pt minus 2pt}
\setlength{\dblfloatsep}{10pt plus 2pt minus 2pt}
\renewcommand{\dbltopfraction}{0.9}
\renewcommand{\dblfloatpagefraction}{0.75}
% Allow for the bullet and gap in the template's abstract/index-term boxes.
% The class aliases these two lengths; separate them before adjusting the box.
\newlength{\AccessAbstractWidth}
\setlength{\AccessAbstractWidth}{\titlewidth}
\addtolength{\AccessAbstractWidth}{-10pt}
\let\abstractwidth\AccessAbstractWidth
% BEGIN EDITABLE PLOT STYLE
% Shared plotting style; all text inherits the article font.
\providecommand{\ArrivalCellFont}{\fontsize{6}{7}\selectfont}
\pgfplotsset{compat=1.18,paper axis/.style={font=\scriptsize,label style={font=\scriptsize},title style={font=\small},tick label style={font=\scriptsize},legend style={font=\scriptsize,draw=none,fill=none},axis line style={black!65},tick style={black!65},grid=major,grid style={black!10},scaled ticks=false,unbounded coords=jump,clip=true}}
\definecolor{FullBlue}{HTML}{2878B5}
\definecolor{ScaleOrange}{HTML}{D45A44}
\definecolor{Size50}{HTML}{D73027}
\definecolor{Size100}{HTML}{FC8D59}
\definecolor{Size250}{HTML}{4575B4}
\definecolor{Size1000}{HTML}{313695}
\pgfplotsset{colormap={paperdiv}{rgb=(0.01960784,0.18823529,0.38039216);
rgb=(0.12725875,0.39584775,0.66874279);
rgb=(0.26274510,0.57647059,0.76470588);
rgb=(0.56647443,0.76870434,0.86851211);
rgb=(0.81960784,0.89803922,0.94117647);
rgb=(0.96908881,0.96647443,0.96493656);
rgb=(0.99215686,0.85882353,0.78039216);
rgb=(0.95455594,0.64175317,0.50572857);
rgb=(0.83921569,0.37647059,0.30196078);
rgb=(0.69227220,0.09227220,0.16770473);
rgb=(0.40392157,0.00000000,0.12156863);}}
\pgfplotsset{colormap={paperseq}{rgb=(1.00000000,1.00000000,0.89803922);
rgb=(1.00000000,0.97539408,0.77193387);
rgb=(0.99764706,0.92156863,0.63607843);
rgb=(0.99607843,0.84347559,0.46915802);
rgb=(0.99607843,0.73490196,0.28000000);
rgb=(0.99497116,0.59747789,0.15949250);
rgb=(0.93960784,0.47137255,0.09490196);
rgb=(0.84822760,0.35229527,0.03497116);
rgb=(0.72000000,0.26039216,0.01098039);
rgb=(0.55686275,0.19123414,0.01737793);
rgb=(0.40000000,0.14509804,0.02352941);}}
% Stream all ECDF vertices to PGF without retaining PGFPlots point metadata.
% Arguments: display-coordinate file, stroke colour, dash setting.
\newcommand{\DrawExactECDF}[3]{%
  \begin{pgfonlayer}{main}%
  \begin{pgfscope}%
    \pgftransformshift{\pgfplotsqpointdescriptionxy{0}{0}}%
    \pgfsetxvec{\pgfpoint{\pgfkeysvalueof{/pgfplots/width}}{0pt}}%
    \pgfsetyvec{\pgfpoint{0pt}{\pgfkeysvalueof{/pgfplots/height}}}%
    \pgfpathrectangle{\pgfpointorigin}{\pgfpoint{\pgfkeysvalueof{/pgfplots/width}}{\pgfkeysvalueof{/pgfplots/height}}}%
    \pgfusepath{clip}%
    \pgfsetstrokecolor{#2}%
    \pgfsetlinewidth{.8pt}%
    #3%
    \pgfplothandlerconstantlineto%
    \pgfplotxyfile{#1}%
    \pgfusepath{stroke}%
  \end{pgfscope}%
  \end{pgfonlayer}%
}
% END EDITABLE PLOT STYLE

\begin{document}
\vol{XX}\AccessYear{XXXX}
\markboth{Appendix A layout check}{Appendix A layout check}
\appendices
\section{\yulian{Verification of Full Gillespie on cliques}{Clique formulas and verification}}
\label{app:clique}

\yulian{We check the Full Gillespie implementation on clique networks, for which mean infection times can be calculated analytically. A clique is a community in which every pair of nodes is connected. We use the first community of the clique chains: it has $n$ nodes and $k$ contacts with the next community. The $k$ source nodes with these contacts are the bridge nodes. The initially infected node and the set of bridge nodes are selected uniformly and independently.

Two times are distinguished. The bridge infection time, $T_B$, is when the first bridge node becomes infected. The export time, $T_{\mathrm{export}}$, is when infection is first transmitted to the next community. Infection of a bridge makes transmission possible, but transmission need not occur immediately. The last column of Table~\ref{tab:verification} measures the additional waiting time, $T_{\mathrm{export}}-T_B$, in Full Gillespie.}{}

\yulian{Table~\ref{tab:verification} compares analytical means with simulation results. For example, for $n=100$ and $k=1$, the simulated mean bridge infection time is $17.05\pm0.38$, compared with the exact value $17.26$. The simulated mean wait from bridge infection to export is $326.5\pm17.0$, compared with the exact value $333.3$. These comparisons check internal spread and, for one contact, the subsequent transmission to the next community.}{}

\begin{table}[!htbp]
\centering
\caption{\yulian{Checks on the first community of the clique chains ($\beta=0.003$, 400 runs per configuration). $T_B$ is the first bridge infection time; $T_{\mathrm{export}}$ is the first transmission time to the next community. Full values are means $\pm$ standard errors. Dashes indicate that the one-contact waiting-time check is not reported for $k=5$.}{Verification of the reference simulator on the first community of the clique chains ($\beta=0.003$, 400 runs per cell): mean bridge-activation time and, for $k=1$, mean wait after activation (exact value $1/\beta=333.3$). Values are mean $\pm$ standard error.}}
\label{tab:verification}
\small
\setlength{\tabcolsep}{3pt}
\begin{tabular}{rrrrr}
\toprule
$n$ & $k$ & \yulian{\shortstack{Exact mean\\$T_B$}}{Exact $\mathbb{E}[T_B]$} & \yulian{\shortstack{\Full{} mean\\$T_B$}}{\Full{} $T_B$} & \yulian{\shortstack{\Full{} mean wait\\$T_{\mathrm{export}}-T_B$}}{\Full{} wait}\\
\midrule
50 & 1 & 29.86 & $29.72\pm0.71$ & $354.0\pm16.9$\\
100 & 1 & 17.26 & $17.05\pm0.38$ & $326.5\pm17.0$\\
250 & 1 & 8.13 & $7.84\pm0.15$ & $347.4\pm16.6$\\
1000 & 1 & 2.49 & $2.51\pm0.04$ & $353.8\pm17.9$\\
50 & 5 & 15.97 & $15.96\pm0.55$ & --\\
100 & 5 & 10.31 & $9.64\pm0.29$ & --\\
250 & 5 & 5.35 & $5.16\pm0.12$ & --\\
1000 & 5 & 1.80 & $1.77\pm0.03$ & --\\
\bottomrule
\end{tabular}
\end{table}

\yulian{For $k=1$, the delay after bridge infection is an independent $\Exp(\beta)$ waiting time, with mean $1/\beta=333.3$ at $\beta=0.003$. For $k=5$, other bridge nodes may become infected before the first export, increasing the total transmission rate. The same constant-rate waiting-time check therefore does not apply. The waiting time is still defined, but this table reports it only for $k=1$; the dashes do not indicate missing runs.}{where $H_{n-1}$ is the harmonic number; the case of a bridge seed ($T_B=0$) is included. For $k=1$, the export time in \Full{} is $T_B$ plus an independent $\Exp(\beta)$ waiting time. The expected time to infect the whole clique is $2H_{n-1}/(\beta n)$. Table~\ref{tab:verification} compares~\eqref{eq:clique_TB} with the first community of the clique chains.}

\yulian{The analytical bridge infection time follows from the SI process on a clique. With $m$ infected and $n-m$ susceptible nodes, there are $m(n-m)$ possible transmitting contacts, each with rate $\beta$. The mean wait for the next internal infection is therefore $1/[\beta m(n-m)]$. The probability that no bridge is among the first $m$ infected nodes is $\binom{n-k}{m}/\binom{n}{m}$. Adding the waiting times weighted by these probabilities gives}{Consider an $n$-node clique with SI dynamics, per-edge rate $\beta$, a uniformly placed seed and a uniformly placed set of $k$ bridge nodes. When $m$ nodes are infected, the next internal infection occurs at rate $\beta m(n-m)$, independently of which nodes are infected, and the infected set is a uniform $m$-subset. The probability that none of the first $m$ infected nodes is a bridge is $\binom{n-k}{m}/\binom{n}{m}$, hence}
\begin{equation}
\label{eq:clique_TB}
\mathbb{E}[T_B]=\sum_{m=1}^{n-k}\frac{\binom{n-k}{m}}{\binom{n}{m}}\frac{1}{\beta m(n-m)},
\end{equation}
\yulian{The case in which the initial infected node is already a bridge has $T_B=0$ and is included in this expectation. For one bridge, the expression simplifies to}{For one bridge this simplifies to}
\begin{equation}
\mathbb{E}[T_B]=\frac{H_{n-1}}{\beta n}\ \ (k=1),
\end{equation}
\yulian{where $H_{n-1}=\sum_{r=1}^{n-1}1/r$ is the harmonic number. The mean time to infect the entire clique is $2H_{n-1}/(\beta n)$.}{}


\EOD\end{document}
